% ============================================================
% Section 3: Physical Background
% BearMamba-3: Mamba-3 + L_kin + Multi-Sensor Fusion
% Target: MSSP (IF 8.4)  Draft: 2026-06-15
% ============================================================

\section{Physical Background}
\label{sec:physics}

\subsection{Bearing Fault Characteristic Frequencies}
\label{subsec:phys_freqs}

Rolling element bearings generate periodic vibration signatures whose
frequencies are analytically determined by bearing geometry and shaft speed.
For a bearing with $n_b$ rolling elements, ball diameter $d$, pitch
diameter $D$, and contact angle $\alpha$, the four primary fault
characteristic frequencies are~\cite{randall2011rolling,smith2015rolling}:
\begin{align}
  f_{\text{BPFO}} &= \frac{n_b}{2}\,f_r\!\left(1 - \frac{d}{D}\cos\alpha\right),
  \label{eq:bpfo} \\
  f_{\text{BPFI}} &= \frac{n_b}{2}\,f_r\!\left(1 + \frac{d}{D}\cos\alpha\right),
  \label{eq:bpfi} \\
  f_{\text{BSF}}  &= \frac{D}{2d}\,f_r\!\left[1 - \left(\frac{d}{D}\cos\alpha\right)^2\right],
  \label{eq:bsf}  \\
  f_{\text{FTF}}  &= \frac{f_r}{2}\!\left(1 - \frac{d}{D}\cos\alpha\right),
  \label{eq:ftf}
\end{align}
where $f_r = \Omega/60$ is the shaft rotation frequency (Hz) and $\Omega$
is shaft speed (rpm).
$f_{\text{BPFO}}$ (Ball Pass Frequency, Outer Race) is excited by outer-race
defects; $f_{\text{BPFI}}$ by inner-race defects; $f_{\text{BSF}}$ (Ball
Spin Frequency) by rolling element defects; and $f_{\text{FTF}}$ (Fundamental
Train Frequency) by cage defects.
Each characteristic frequency and its integer harmonics produce periodic
impulse trains in the measured vibration signal, which form the primary
diagnostic signatures for bearing condition monitoring and fault
detection~\cite{lei2020applications,antoni2006spectral_kurtosis}.

Table~\ref{tab:bearing_params} lists the bearing parameters used in this
study and the resulting characteristic frequencies at nominal shaft speeds.

\begin{table}[htbp]
\centering
\caption{Bearing parameters and fault characteristic frequencies at
         nominal shaft speed for each dataset.}
\label{tab:bearing_params}
\begin{tabular}{lcccccc}
\toprule
Dataset & Bearing & $n_b$ & $d$ (mm) & $D$ (mm) & $\Omega$ (rpm)
        & $f_{\text{BPFO}}$ / $f_{\text{BPFI}}$ (Hz) \\
\midrule
CWRU DE & SKF 6205 & 9 & 7.94  & 39.04 & 1797
         & 107.36 / 162.19 \\
CWRU FE & SKF 6203 & 8 & 6.75  & 28.50 & 1797
         & 91.44$^\dagger$ / 148.16$^\dagger$ \\
PU      & SKF 6203 & 8 & 6.75  & 29.05 & 1500
         & 76.77 / 123.23 \\
XJTU-SY & LDK UER204 & 8 & 7.92 & 34.55 & 2400
          & 123.32 / 196.68 \\
\bottomrule
\end{tabular}
\begin{tablenotes}
\item[$\dagger$] CWRU FE bearing (SKF 6203) has a different pitch diameter
  (28.50~mm) from the PU bearing of the same part number (29.05~mm);
  these are computed separately.
  All contact angles $\alpha = 0^\circ$.
\end{tablenotes}
\end{table}

\subsection{Speed-Dependent Shift and Per-Sample Frequency Targets}
\label{subsec:phys_speed}

Because all characteristic frequencies scale proportionally with shaft
speed $f_r$, any variation in operating speed directly shifts the
entire fault frequency spectrum.
Under variable-speed conditions — such as the accelerated-life XJTU-SY
recordings where instantaneous speed fluctuates within a measurement window
— using a fixed dataset-level frequency set would introduce misalignment
errors proportional to the speed deviation~\cite{liu2023var_speed,liao2020variable_domain}.

This constraint motivates computing the fault frequency target set
$\mathcal{F}_i$ \emph{per sample} from each sample's known shaft speed
$\Omega_i$, which is available as metadata in all three datasets used here
(CWRU nominal 1797~rpm; PU nominal 900 or 1500~rpm; XJTU-SY nominal 2100,
2250, or 2400~rpm per condition).
This per-sample design makes $\mathcal{L}_{\text{kin}}$ naturally applicable
to variable-speed scenarios without modification to the loss function —
a key distinction from fixed-spectrum regularisers.

Spectral Kurtosis~\cite{antoni2006spectral_kurtosis} represents the
classical signal-processing approach for identifying impulsive
fault-frequency content; envelope demodulation then localises BPFO/BPFI
peaks.
BearMamba-3 extends this principle to the learned-representation domain
by encouraging the Mamba-3 backbone's internal oscillatory states to
resonate at the analytically known fault harmonics during training
(Section~\ref{subsec:lkin}), while preserving end-to-end trainability
from raw waveforms.
